\documentclass[noanswers, addpoints, 12pt]{exam}
\usepackage[T1]{fontenc}
\usepackage[utf8]{inputenc}
\usepackage{lmodern}
\usepackage[french]{babel}
\usepackage{xcolor}
\usepackage{listings}
\usepackage{graphicx}
\usepackage[most]{tcolorbox}

\definecolor{codegreen}{rgb}{0,0.6,0}
\definecolor{codegray}{rgb}{0.5,0.5,0.5}
\definecolor{codepurple}{rgb}{0.58,0,0.82}
\definecolor{backcolour}{rgb}{0.95,0.95,0.92}
 
\lstdefinestyle{mystyle}{
    backgroundcolor=\color{backcolour},   
    commentstyle=\color{codegreen},
    keywordstyle=\color{magenta},
    numberstyle=\tiny\color{codegray},
    stringstyle=\color{codepurple},
    basicstyle=\footnotesize,
    breakatwhitespace=false,         
    breaklines=true,                 
    captionpos=b,                    
    keepspaces=true,                 
    numbers=left,                    
    numbersep=5pt,                  
    showspaces=false,                
    showstringspaces=false,
    showtabs=false,                  
    tabsize=2
}
\lstset{style=mystyle, language=Python}

\title{CC1 Prem NSI - Conditions et boucles (rattrapage)}
\author{Babash - NSI Tani Malandi et Johanne Bakalara}
\date{Octobre 2024, Durée: 1h}


\lstset{extendedchars=true,literate={é}{{\'e}}1}
\begin{document}
\runningfooter{}{\footnotesize Babash - NSI Tani Malandi et Johanne Bakalara \textperiodcentered{} CC BY-SA 4.0 \textperiodcentered{} page \thepage}{}
\firstpagefooter{}{\footnotesize Babash - NSI Tani Malandi et Johanne Bakalara \textperiodcentered{} CC BY-SA 4.0 \textperiodcentered{} page \thepage}{}

\maketitle

\pagestyle{headandfoot}
\thispagestyle{headandfoot}
\firstpageheader{Nom : \\  Prénom :\\ Classe :}{}{}


\section*{Consignes:}
\begin{center} 
\fbox{\parbox{\linewidth}{
\begin{center}
Le devoir est à rédiger sur une feuille libre.
% La clarté de la rédaction et les explications font partie intégrante de la notation. Une réponse non rédigée ou sans explication sera considérée fausse.\\
Prenez le temps d'encadrer les réponses.
\end{center}
}}
\end{center}

\vsword{Note}
\hsword{Note}

\section*{Barème :}
\begin{center}
  \par\bigskip
  \gradetable[h]
\end{center}
\begin{questions}

\section{Question de cours : }

\question[2]  :
\begin{parts}
    \part Quelles sont les 5 structures de contrôle vu jusqu'ici ?
    \part Quelle fonction permet de récupérer la saisie d'un utilisateur en Python ?
    \part Quelle fonction permet d'afficher une information en Python ?
    \part Que signifie le terme indentation ?
\end{parts}

\question[2] Donner les plages de valeurs correspondant aux appels de range suivants :
\begin{itemize}
    \item range(3)
    \item range(9, 14)
    \item range(5, 15, 5)
    \item range(64, 12, -10)
\end{itemize}

\question[2] Quels opérateurs Python permettent de réaliser les opérations suivantes :
\begin{itemize}
    \item attribuer une valeur à une variable :
    \item calculer le quotient d'une division entiere :
    \item tester si deux expressions sont égales :
    \item faire la somme de deux expressions :
\end{itemize}
 
\section{Variables et conditions : }

\begin{lstlisting}[language=Python]
    import random
    # random.randint(1, 10) tire un nombre aléatoire entre 1 et 10.
    pari = int(input("Saisissez un nombre entre 1 et 36 : "))
    score = 0
    roulette = random.randint(1, 36)
    if  roulette == pari:
        score = score + 35
    else :
        score = 0
    if -4 <= roulette - pari <= 4 :
        score = score + 20
    if (roulette % 2) == (pari % 2):
        score = score + 2
    elif score != 0 :
        score = score - 2 
\end{lstlisting}
\question[1] Combien y a-t-il de variables dans ce code, quelles sont leur type ?
\question[3] Donner le tableau d'avancement des variables pour pari=10, roulette=13.
\question[2] Quel est le score pour :
\begin{itemize}
    \item pari=2, roulette=2
    \item pari=10, roulette=25
    \item pari=9, roulette=33
    \item pari=5, roulette=7
\end{itemize}

\section{Boucles : }

\question[1] Combien y a-t-il de variables dans ce code, quelles sont leur type ?
\question[4] Faire le tableau d'avancement des variables du code ci-dessous :
\begin{lstlisting}[language=Python]
n = 12
i = 0
p = 0
while (i*p) < n:
    p = p + 1
    i = i + 1
\end{lstlisting}

\section{Rédaction de code : }
\question[3] Écrire un code permettant de calculer la somme suivante : $$\frac{1}{2} + \frac{2}{3} +\frac{3}{4} + ... + \frac{99}{100}$$ Vous devez impérativement utiliser une boucle.


\end{questions}

\end{document}

